\chapter{Instalação, operação, testes e manutenção}\label{cap:operacao}
\section{Preparação do ambiente e primeira partida}
Instalar Python e PlatformIO com suporte a ESP-IDF, obter a revisão correta do repositório e registrar o commit. Conferir \texttt{platformio.ini} de cada alvo e as partições OTA. Os painéis web são gerados a partir das fontes HTML/CSS/JavaScript por \texttt{gerar\_painel.py}; o arquivo C incorporado deve corresponder a elas. Credenciais Wi-Fi, URL e token do Google pertencem à configuração local ignorada pelo Git. Não colocar segredos no ZIP do Overleaf nem no histórico público.

Na primeira instalação, gravar cada placa por cabo para criar a tabela de partições. Os comandos e ambientes corretos estão no capítulo~\ref{cap:ota}; selecionar a porta e identificar fisicamente cada placa antes da gravação.
O ambiente exato do mestre fixa \texttt{espressif32@6.13.0}; os satélites devem ser compilados com dependências registradas no ensaio, pois seus arquivos não fixam essa versão. O terminal serial do mestre e a chave inicial da equipe usam 115\,200 baud. A gravação por cabo não dispensa conferir GPIO, alimentação e conexão física antes de ligar o conjunto.

\subsection{Montagem e partida}
\begin{enumerate}
\item Identificar a revisão física de cada PCB e o esquema correspondente; verificar polaridade, tensão dos conversores, proteções, continuidade de GND e ausência de curto com a alimentação desligada.
\item Confirmar 3,3\,V nos domínios lógicos. Conferir pull-ups dos quatro canais de encoder para 3,3\,V, nunca para a alimentação externa. Conferir TX cruzado com RX em cada UART.
\item Confirmar geometria das rodas e relação encoder/roda; conferir orientação X/Y do MPU; registrar valores efetivos antes de análise quantitativa.
\item Energizar, observar serial, LEDs e OLED. Verificar se o mestre inicia, se o sensor inercial calibra em repouso, se os encoders não contam com rodas paradas e se o cartão monta quando instalado.
\item Acessar o AP da Equipe em \url{http://192.168.4.1/}; conferir versão e idade da telemetria dos dois satélites, sequência e ACK. Verificar se a origem indicada é real ou simulada.
\item Iniciar um registro curto, conferir o contador de amostras no SD e preservar o CSV original. Para ensaio do painel público, verificar que Visitantes mostra somente leitura.
\end{enumerate}

\subsection{Conferência elétrica antes da primeira energização}\label{sec:conferencia-eletrica}
Trabalhar com a alimentação de tração isolada e descarregada segundo o procedimento elétrico da equipe. Identificar a revisão efetivamente montada; não usar apenas os nomes dos conectores da folha v4. Com multímetro, confirmar polaridade dos trilhos, continuidade dos retornos, ausência de curto entre 48 V, 5 V, 3,3 V e GND, e orientação do diodo de roda livre. Energizar a eletrônica por fonte limitada em corrente antes de conectá-la ao veículo. Medir 5 V e 3,3 V em cada módulo durante partida, Wi-Fi e escrita no cartão, além da temperatura do buck e do regulador. Registrar queda de tensão e ruído com osciloscópio quando possível. O capítulo \ref{cap:eletronica} explica as razões elétricas e os pontos não comprovados do esquema.

Para o I\textsuperscript{2}C, contar todos os resistores de pull-up soldados na placa e nos módulos. Medir SDA--3,3 V e SCL--3,3 V com o conjunto desligado, comparar a resistência equivalente com capacitância e tempo de subida, e retirar os pull-ups excedentes dos módulos quando necessário. Conferir que nenhum módulo leva as linhas a 5 V. Observar uma transação a 400 kHz com todos os módulos conectados; registrar nível alto, nível baixo, subida e ausência de colisões. Só então verificar endereços de MPU, OLED e INA228, caso este último esteja instalado. O driver do INA228 está implementado, mas sua resposta depende de módulo presente, ligação correta e configuração coerente com o shunt físico.

Para o shunt, confirmar que o componente real é de 1 m$\Omega$, sua potência admissível contínua e a montagem térmica a 75 A. Usar ligação Kelvin independente dos cabos de potência, marcar IN+ e IN$-$ e comparar a tensão diferencial com $1\,\mathrm{mV/A}$ em correntes conhecidas e seguras. Medir aquecimento em regime; 75 A produzem 5,625 W no elemento resistivo. Não usar leitura de corrente, potência ou estado de bateria como proteção até que os limites e os ensaios da montagem sejam aceitos; o driver de leitura, sozinho, não implementa proteção de tração.

O relé e seu transistor devem ser ensaiados separados do controlador de tração: identificar tensão nominal da bobina, corrente com o contato aberto e fechado, estado em reset e falta de alimentação, polaridade de D1, tempo de soltura e contato normalmente aberto/fechado. A topologia representada na folha v4 exige revisão antes de qualquer conexão ao driver. O estado remoto de parada descrito no capítulo \ref{cap:comandos} não aciona esse relé.

\subsection{Configuração de rede, painéis e registro}\label{sec:config-operacional}
Depois de gravar os três ESPs, observar no serial se as partições e a NVS foram inicializadas sem erro. Conectar o computador ou telefone ao AP da Equipe e abrir \url{http://192.168.4.1/}; registrar nome da rede, canal, versão, número de clientes e idade da telemetria. O acesso ao AP Equipe permite enviar comandos sem chave HTTP adicional; reservar a senha desse AP à equipe. Conferir a página de Visitantes separadamente para garantir que ela apresenta dados e não oferece comandos. O capítulo \ref{cap:wifi-paineis} descreve as rotas, a política de clientes e o motivo de o mestre ligar o Wi-Fi externo somente durante sessões.

Na Equipe ou no terminal do mestre, inspecionar perfis Wi-Fi salvos, tentar a rede principal e registrar canal, RSSI, tentativas, erro e prazo. Para uma rede secundária, registrar autorização e a política de conexão automática antes de testar. Uma conexão ao roteador não comprova acesso ao Google nem integridade do relógio ou da URL HTTPS. Fazer uma consulta OTA e um envio de CSV de ensaio separadamente, observando o estado final de cada tarefa. Para o Google, seguir o cadastro de propriedades, implantação e teste do Apps Script do capítulo \ref{cap:memorias_google}; não copiar token, chave ou SSID para este manual.

\subsection{Ensaios de falha e recuperação}\label{sec:ensaios-falha}
Verificar deliberadamente, em bancada, situações que o diagnóstico deve distinguir: sensor I\textsuperscript{2}C ausente, pull-up inadequado, UART de um satélite interrompida, Wi-Fi externo indisponível, cartão SD removido antes da partida, memória sob carga de TLS, resposta HTTP de upload inválida e imagem OTA de teste que falhe a validação. Para cada caso, registrar causa ativa, contador, página OLED/LED, resposta web, log, tempo até detecção, comportamento preservado e recuperação. Repetir após religar a alimentação; a memória RTC usada pelo guardião de OTA não é persistência garantida sob corte total. Não provocar falha por desconexão de cabos energizados de potência ou por remoção do cartão durante escrita.

Ao testar rollback, gravar uma imagem de laboratório que reinicie antes de ser validada e confirmar qual partição foi selecionada no próximo boot. Em teste separado, provocar três reinícios anormais da mesma imagem após a confirmação inicial e dentro da janela de estabilização; observar contagem e retorno. Esses são dois mecanismos diferentes, descritos no capítulo \ref{cap:ota}; registrar sequência temporal, causa de reset e versão efetiva em cada reinício. Não realizar o ensaio com controle de tração conectado.

\section{Operação e manutenção}
Antes de mover o veículo, conferir alimentação, conectores, fixação dos sensores, cabos afastados da potência, estado dos LEDs e integridade da parada física do veículo conforme o procedimento próprio da equipe. O painel remoto ainda não é mecanismo de parada física. Deixar o MPU em repouso durante a calibração inicial. Confirmar validade, origem e idade da telemetria. Ao encerrar, guardar os CSVs e logs do ensaio junto da versão de firmware e das condições experimentais. Desligar conforme a sequência elétrica adotada na montagem, evitando remover o cartão durante escrita.

\subsection{Plano de verificação}
Os critérios quantitativos de aceitação devem ser aprovados pela equipe para a finalidade do veículo; não há no material consultado uma campanha de calibração conclusiva. A tabela descreve ensaios reproduzíveis e a evidência a guardar.
\begin{table}[htbp]\centering\small
\caption{Ensaios recomendados antes de declarar uma medição validada}\label{tab:ensaios}
\begin{tabularx}{\textwidth}{@{}l X X@{}}\toprule
Objeto & Procedimento & Evidência e decisão\\\midrule
Alimentação & Medir trilhos em vazio, partida e carga; inspecionar ruído e aquecimento. & Tensão, instrumento, corrente, temperatura e comparação com limites dos componentes.\\
Encoders & Girar cada roda número conhecido de voltas, nos dois sentidos e com carro apoiado; comparar distância externa. & Contagens, $D$, $i$, erro e repetibilidade por lado; confirmar ausência de sentido na versão atual.\\
MPU6050 & Calibrar parado; verificar eixos por inclinação e movimento controlado. & Estado de calibração, orientação, ruído e resposta longitudinal.\\
UART & Interromper e restaurar cada enlace; injetar quadro corrompido em bancada. & Contadores CRC/perdas/duplicatas, idade da amostra e recuperação independente.\\
SD & Gravar por tempo suficiente, retirar alimentação de modo controlado e reiniciar. & CSV legível, contadores, rotação, perda máxima observada e preservação dos arquivos.\\
OTA & Instalar versão de teste, verificar alvo, reinício, versão nova e rollback controlado. & Manifesto, SHA-256, logs, estado após observação e resultado por dispositivo.\\\bottomrule
\end{tabularx}\end{table}

Não classificar o ensaio apenas como aprovado ou reprovado. Registrar data, responsável, revisão da placa, commit, versões dos três firmwares, instrumentos, condições, dados brutos, diferença em relação à referência e critério previamente definido. Testar também o sistema com Wi-Fi externo indisponível, SD ausente e navegador desconectado: a aquisição e a telemetria UART devem continuar quando a falha for degradável.

\subsection{Atualização de firmware}
Para OTA, conferir se há release completo com manifesto e três binários coerentes. Consultar a versão disponível no menu 2 do mestre ou na página Comandos da Equipe, autorizar o alvo, acompanhar progresso e aguardar a versão nova responder e completar a observação. Uma oferta sem autorização expira; verificar novamente se necessário. Em atualização de todos, satélites vêm antes do mestre. Não interpretar confirmação HTTP como instalação concluída.

\begin{table}[htbp]\centering\small
\caption{Sintomas e primeiras verificações}\label{tab:diagnostico}
\begin{tabularx}{\textwidth}{@{}l X X@{}}\toprule
Sintoma & Verificar & Ação inicial\\\midrule
Sem amostra recente & Alimentação dos três ESPs, TX/RX/GND, baud, CRC, sequência e ACK. & Restaurar o enlace, consultar contadores e não confiar no último valor exibido.\\
Contagem com roda parada & Pull-ups a 3,3\,V, blindagem, ruído de potência e filtro PCNT. & Isolar o chicote e medir formas de onda antes de mudar a geometria.\\
Velocidade incoerente & Diâmetro, relação de transmissão, P/R e quatro bordas contadas. & Comparar contagens por volta e aferir a distância com referência.\\
MPU inválido & SDA/SCL, endereço, alimentação, eixo e condição de repouso. & Conferir I\textsuperscript{2}C e repetir calibração após fixar a orientação.\\
Sem CSV & Cartão, 3,3\,V, SPI2, montagem FAT, fila e espaço livre. & Preservar o cartão e consultar logs; telemetria ao vivo pode seguir.\\
Ação web negada & Permissão do Mestre (quando exigida), carga inválida ou idade da UART. & Conferir a mensagem de rejeição, a permissão aplicável e consultar resposta e estado separadamente.\\
OTA interrompida & Rede externa, release, manifesto, hash, espaço da partição e ACK UART. & Ler erro por alvo e versão que responde antes de repetir.\\\bottomrule
\end{tabularx}\end{table}

\subsection{Pendências técnicas rastreáveis}
\begin{enumerate}
\item Associar a placa física instalada a um projeto KiCad, revisão de esquema, PCB, conectores e mapa de alimentação; acrescentar fotos de ambas as faces e do chicote.
\item Medir diâmetro efetivo das rodas, relação real dos encoders, resolução e erro; validar orientação do MPU6050 no chassi.
\item Integrar e calibrar o INA228 e o shunt para corrente e tensão do barramento; implementar os demais sensores físicos antes de publicar seus campos como medições.
\item Registrar testes de desempenho com SD, Wi-Fi e OTA concorrentes e os prazos de detecção de amostra antiga; aprovar limites de erro e critérios por aplicação.
\item Se houver futura atuação física, projetar uma cadeia independente de segurança, intertravamentos, autenticação de operador, prazo de validade dos pedidos e retorno de estado medido. O estado preparatório atual não atende essa função.
\end{enumerate}
